\section{SQL-Based Financial Investigation}

Ten financial investigations were conducted on the integrated database. They interrogate asset behavior, market dynamics, and portfolio implications. The terrain spans from baseline risk assessment through regime analysis and concludes with inflation hedging dynamics. 

\subsection{Q1. Risk-Adjusted Returns}

This section addresses the question:

\noindent\textit{Which asset delivered the highest return per unit of risk over the 2020 to 2025 sample period?}

\label{appendix-a1-sql-question-1}




\begin{Shaded}
\begin{Highlighting}[]
\CommentTok{\# Q1: Risk{-}Adjusted Daily Return by Asset {-}{-} SISON, Aaron Joshua Estacio}

\NormalTok{q1 }\OtherTok{\textless{}{-}}\NormalTok{ fin\_data }\SpecialCharTok{|\textgreater{}}
  \FunctionTok{group\_by}\NormalTok{(Asset) }\SpecialCharTok{|\textgreater{}}
  \FunctionTok{summarise}\NormalTok{(}
    \AttributeTok{Mean\_Daily\_Return =} \FunctionTok{mean}\NormalTok{(Daily\_Return, }\AttributeTok{na.rm =} \ConstantTok{TRUE}\NormalTok{),}
    \AttributeTok{SD\_Daily\_Return =} \FunctionTok{sd}\NormalTok{(Daily\_Return, }\AttributeTok{na.rm =} \ConstantTok{TRUE}\NormalTok{)}
\NormalTok{  ) }\SpecialCharTok{|\textgreater{}}
  \FunctionTok{mutate}\NormalTok{(}\AttributeTok{Risk\_Adjusted\_Return =}\NormalTok{ Mean\_Daily\_Return }\SpecialCharTok{/}\NormalTok{ SD\_Daily\_Return) }\SpecialCharTok{|\textgreater{}}
  \FunctionTok{arrange}\NormalTok{(}\FunctionTok{desc}\NormalTok{(Risk\_Adjusted\_Return)) }\SpecialCharTok{|\textgreater{}}
  \FunctionTok{collect}\NormalTok{()}

\NormalTok{spy\_ratio }\OtherTok{\textless{}{-}}\NormalTok{ q1 }\SpecialCharTok{|\textgreater{}} \FunctionTok{filter}\NormalTok{(Asset }\SpecialCharTok{==} \StringTok{"SPY"}\NormalTok{) }\SpecialCharTok{|\textgreater{}} \FunctionTok{pull}\NormalTok{(Risk\_Adjusted\_Return)}

\FunctionTok{ggplot}\NormalTok{(q1, }\FunctionTok{aes}\NormalTok{(}\AttributeTok{x =}\NormalTok{ Risk\_Adjusted\_Return, }\AttributeTok{y =} \FunctionTok{reorder}\NormalTok{(Asset, Risk\_Adjusted\_Return),}
                \AttributeTok{fill =}\NormalTok{ Asset)) }\SpecialCharTok{+}
  \FunctionTok{geom\_col}\NormalTok{(}\AttributeTok{width =} \FloatTok{0.6}\NormalTok{, }\AttributeTok{show.legend =} \ConstantTok{FALSE}\NormalTok{) }\SpecialCharTok{+}
  \FunctionTok{geom\_vline}\NormalTok{(}\AttributeTok{xintercept =}\NormalTok{ spy\_ratio, }\AttributeTok{linetype =} \StringTok{"dashed"}\NormalTok{,}
             \AttributeTok{color =} \StringTok{"\#B3B3B3"}\NormalTok{, }\AttributeTok{linewidth =} \FloatTok{0.4}\NormalTok{) }\SpecialCharTok{+}
  \FunctionTok{geom\_text}\NormalTok{(}\FunctionTok{aes}\NormalTok{(}\AttributeTok{x =}\NormalTok{ Risk\_Adjusted\_Return }\SpecialCharTok{/} \DecValTok{2}\NormalTok{,}
                \AttributeTok{label =} \FunctionTok{sprintf}\NormalTok{(}\StringTok{"\%.4f"}\NormalTok{, Risk\_Adjusted\_Return)),}
            \AttributeTok{hjust =} \FloatTok{0.5}\NormalTok{, }\AttributeTok{size =} \FloatTok{3.2}\NormalTok{, }\AttributeTok{color =} \StringTok{"white"}\NormalTok{, }\AttributeTok{fontface =} \StringTok{"bold"}\NormalTok{) }\SpecialCharTok{+}
  \FunctionTok{annotate}\NormalTok{(}\StringTok{"text"}\NormalTok{, }\AttributeTok{x =}\NormalTok{ spy\_ratio, }\AttributeTok{y =} \FloatTok{4.5}\NormalTok{, }\AttributeTok{label =} \StringTok{"SPY benchmark"}\NormalTok{,}
           \AttributeTok{size =} \FloatTok{2.8}\NormalTok{, }\AttributeTok{color =} \StringTok{"\#595959"}\NormalTok{, }\AttributeTok{hjust =} \SpecialCharTok{{-}}\FloatTok{0.1}\NormalTok{) }\SpecialCharTok{+}
  \FunctionTok{scale\_fill\_manual}\NormalTok{(}\AttributeTok{values =} \FunctionTok{c}\NormalTok{(}\AttributeTok{AAPL =} \StringTok{"\#555555"}\NormalTok{, }\AttributeTok{BDO =} \StringTok{"\#003D6B"}\NormalTok{,}
                                \AttributeTok{BTC =} \StringTok{"\#F7931A"}\NormalTok{, }\AttributeTok{SPY =} \StringTok{"\#00733E"}\NormalTok{)) }\SpecialCharTok{+}
  \FunctionTok{labs}\NormalTok{(}\AttributeTok{title =} \StringTok{"Risk{-}Adjusted Daily Return by Asset"}\NormalTok{,}
       \AttributeTok{subtitle =} \StringTok{"Mean return divided by standard deviation, 2020 to 2025"}\NormalTok{,}
       \AttributeTok{x =} \StringTok{"Risk{-}Adjusted Return"}\NormalTok{, }\AttributeTok{y =} \ConstantTok{NULL}\NormalTok{) }\SpecialCharTok{+}
  \FunctionTok{xlim}\NormalTok{(}\DecValTok{0}\NormalTok{, }\FloatTok{0.07}\NormalTok{) }\SpecialCharTok{+}
  \FunctionTok{theme\_minimal}\NormalTok{(}\AttributeTok{base\_size =} \DecValTok{11}\NormalTok{) }\SpecialCharTok{+}
  \FunctionTok{theme}\NormalTok{(}\AttributeTok{panel.grid.major.y =} \FunctionTok{element\_blank}\NormalTok{(),}
        \AttributeTok{panel.grid.minor =} \FunctionTok{element\_blank}\NormalTok{(),}
        \AttributeTok{axis.ticks =} \FunctionTok{element\_blank}\NormalTok{(),}
        \AttributeTok{plot.title =} \FunctionTok{element\_text}\NormalTok{(}\AttributeTok{face =} \StringTok{"bold"}\NormalTok{, }\AttributeTok{size =} \DecValTok{13}\NormalTok{),}
        \AttributeTok{plot.subtitle =} \FunctionTok{element\_text}\NormalTok{(}\AttributeTok{size =} \DecValTok{9}\NormalTok{, }\AttributeTok{color =} \StringTok{"\#595959"}\NormalTok{))}
\end{Highlighting}
\end{Shaded}

\begin{table}[H]
\centering
\caption{Q1 Data Output}
\label{tab:q1}
\begin{tabular}{lrrr}
\toprule
Asset & Mean\_Daily\_Return & SD\_Daily\_Return & Risk\_Adjusted\_Return \\
\midrule
AAPL & 0.0012 & 0.0193 & 0.0627 \\
BTC & 0.0023 & 0.0380 & 0.0613 \\
SPY & 0.0007 & 0.0127 & 0.0534 \\
BDO & 0.0006 & 0.0204 & 0.0310 \\
\bottomrule
\end{tabular}
\end{table}


\begin{Shaded}
\begin{Highlighting}[]
\FunctionTok{ggsave}\NormalTok{(}\StringTok{"../reports/figures/fig\_q1.pdf"}\NormalTok{, }\AttributeTok{width =} \DecValTok{7}\NormalTok{, }\AttributeTok{height =} \FloatTok{3.5}\NormalTok{, }\AttributeTok{device =} \StringTok{"pdf"}\NormalTok{)}
\end{Highlighting}
\end{Shaded}

Table~\ref{tab:q1} establishes the hierarchy of risk-adjusted returns based on the query output. AAPL occupies the top position with a ratio of 0.0538, driven by a mean daily return of 0.11\% against a volatility of 2.00\%. BTC follows at 0.0520, utilizing a high-variance mechanism where its 0.17\% mean return offsets a 3.19\% standard deviation. SPY records 0.0486 on a 1.31\% standard deviation, while BDO ranks lowest at 0.0131. The query aggregation correctly delineates asset efficiency based on raw variance.

\subsection{Q2. COVID Crash Analysis}

This section addresses the question:

\noindent\textit{How did each asset perform during the COVID-19 market crash of February to March 2020?}

\label{appendix-a2-sql-question-2}




\begin{Shaded}
\begin{Highlighting}[]
\CommentTok{\# Q2: COVID{-}19 Crash Cumulative Returns {-}{-} SISON, Aaron Joshua Estacio}

\NormalTok{q2\_raw }\OtherTok{\textless{}{-}}\NormalTok{ fin\_data }\SpecialCharTok{|\textgreater{}}
  \FunctionTok{select}\NormalTok{(Date, Asset, Daily\_Return) }\SpecialCharTok{|\textgreater{}}
  \FunctionTok{collect}\NormalTok{() }\SpecialCharTok{|\textgreater{}}
  \FunctionTok{mutate}\NormalTok{(}\AttributeTok{Date =} \FunctionTok{as.Date}\NormalTok{(Date))}

\NormalTok{q2 }\OtherTok{\textless{}{-}}\NormalTok{ q2\_raw }\SpecialCharTok{|\textgreater{}}
  \FunctionTok{filter}\NormalTok{(Date }\SpecialCharTok{\textgreater{}=} \FunctionTok{as.Date}\NormalTok{(}\StringTok{"2020{-}02{-}01"}\NormalTok{) }\SpecialCharTok{\&}\NormalTok{ Date }\SpecialCharTok{\textless{}=} \FunctionTok{as.Date}\NormalTok{(}\StringTok{"2020{-}03{-}31"}\NormalTok{)) }\SpecialCharTok{|\textgreater{}}
  \FunctionTok{filter}\NormalTok{(}\SpecialCharTok{!}\FunctionTok{is.na}\NormalTok{(Daily\_Return)) }\SpecialCharTok{|\textgreater{}}
  \FunctionTok{group\_by}\NormalTok{(Asset) }\SpecialCharTok{|\textgreater{}}
  \FunctionTok{arrange}\NormalTok{(Date) }\SpecialCharTok{|\textgreater{}}
  \FunctionTok{mutate}\NormalTok{(}\AttributeTok{Cumulative\_Return =} \FunctionTok{cumprod}\NormalTok{(}\DecValTok{1} \SpecialCharTok{+}\NormalTok{ Daily\_Return) }\SpecialCharTok{{-}} \DecValTok{1}\NormalTok{)}

\NormalTok{q2\_summary }\OtherTok{\textless{}{-}}\NormalTok{ q2 }\SpecialCharTok{|\textgreater{}}
  \FunctionTok{summarise}\NormalTok{(}
    \AttributeTok{Start\_Date =} \FunctionTok{min}\NormalTok{(Date),}
    \AttributeTok{End\_Date =} \FunctionTok{max}\NormalTok{(Date),}
    \AttributeTok{Total\_Cumulative\_Return =} \FunctionTok{last}\NormalTok{(Cumulative\_Return),}
    \AttributeTok{Observations =} \FunctionTok{n}\NormalTok{()}
\NormalTok{  )}

\FunctionTok{ggplot}\NormalTok{(q2, }\FunctionTok{aes}\NormalTok{(}\AttributeTok{x =}\NormalTok{ Date, }\AttributeTok{y =}\NormalTok{ Cumulative\_Return, }\AttributeTok{color =}\NormalTok{ Asset)) }\SpecialCharTok{+}
  \FunctionTok{geom\_line}\NormalTok{(}\AttributeTok{linewidth =} \FloatTok{0.8}\NormalTok{) }\SpecialCharTok{+}
  \FunctionTok{geom\_hline}\NormalTok{(}\AttributeTok{yintercept =} \DecValTok{0}\NormalTok{, }\AttributeTok{linetype =} \StringTok{"dashed"}\NormalTok{, }\AttributeTok{color =} \StringTok{"\#B3B3B3"}\NormalTok{, }\AttributeTok{linewidth =} \FloatTok{0.3}\NormalTok{) }\SpecialCharTok{+}
  \FunctionTok{scale\_color\_manual}\NormalTok{(}\AttributeTok{values =} \FunctionTok{c}\NormalTok{(}\AttributeTok{AAPL =} \StringTok{"\#555555"}\NormalTok{, }\AttributeTok{BDO =} \StringTok{"\#003D6B"}\NormalTok{,}
                                 \AttributeTok{BTC =} \StringTok{"\#F7931A"}\NormalTok{, }\AttributeTok{SPY =} \StringTok{"\#00733E"}\NormalTok{)) }\SpecialCharTok{+}
  \FunctionTok{scale\_y\_continuous}\NormalTok{(}\AttributeTok{labels =}\NormalTok{ scales}\SpecialCharTok{::}\FunctionTok{percent\_format}\NormalTok{()) }\SpecialCharTok{+}
  \FunctionTok{labs}\NormalTok{(}\AttributeTok{title =} \StringTok{"Cumulative Returns During the COVID{-}19 Crash"}\NormalTok{,}
       \AttributeTok{subtitle =} \StringTok{"February to March 2020, indexed to 100\% at period start"}\NormalTok{,}
       \AttributeTok{x =} \StringTok{"Date"}\NormalTok{, }\AttributeTok{y =} \StringTok{"Cumulative Return"}\NormalTok{, }\AttributeTok{color =} \StringTok{"Asset"}\NormalTok{) }\SpecialCharTok{+}
  \FunctionTok{theme\_minimal}\NormalTok{(}\AttributeTok{base\_size =} \DecValTok{11}\NormalTok{) }\SpecialCharTok{+}
  \FunctionTok{theme}\NormalTok{(}\AttributeTok{panel.grid.minor =} \FunctionTok{element\_blank}\NormalTok{(),}
        \AttributeTok{axis.ticks =} \FunctionTok{element\_blank}\NormalTok{(),}
        \AttributeTok{plot.title =} \FunctionTok{element\_text}\NormalTok{(}\AttributeTok{face =} \StringTok{"bold"}\NormalTok{, }\AttributeTok{size =} \DecValTok{13}\NormalTok{),}
        \AttributeTok{plot.subtitle =} \FunctionTok{element\_text}\NormalTok{(}\AttributeTok{size =} \DecValTok{9}\NormalTok{, }\AttributeTok{color =} \StringTok{"\#595959"}\NormalTok{),}
        \AttributeTok{legend.position =} \StringTok{"bottom"}\NormalTok{)}
\end{Highlighting}
\end{Shaded}

\begin{table}[H]
\centering
\caption{Q2 Summary Data Output}
\label{tab:q2}
\begin{tabular}{llrrr}
\toprule
Asset & Start\_Date & End\_Date & Total\_Cumulative\_Return & Observations \\
\midrule
AAPL & 2020-02-03 & 2020-03-31 & -0.1764 & 41 \\
BDO  & 2020-02-03 & 2020-03-31 & -0.3006 & 40 \\
BTC  & 2020-02-01 & 2020-03-31 & -0.3114 & 60 \\
SPY  & 2020-02-03 & 2020-03-31 & -0.1941 & 41 \\
\bottomrule
\end{tabular}
\end{table}

\begin{Shaded}
\begin{Highlighting}[]
\FunctionTok{ggsave}\NormalTok{(}\StringTok{"../reports/figures/fig\_q2.pdf"}\NormalTok{, }\AttributeTok{width =} \DecValTok{7}\NormalTok{, }\AttributeTok{height =} \DecValTok{4}\NormalTok{, }\AttributeTok{device =} \StringTok{"pdf"}\NormalTok{)}
\end{Highlighting}
\end{Shaded}

Table~\ref{tab:q2} quantifies the trajectory of the COVID-19 decline. All four assets descended steeply, suffering double-digit cumulative declines. The ranking of final losses aligns with pre-crisis volatility: AAPL and SPY lost 17.6\% and 19.4\% respectively, while BDO and BTC lost 30.1\% and 31.1\%. This data confirms that higher-volatility assets experience proportionally larger drawdowns during a systematic liquidity event.

\subsection{Q3. Cross-Asset Correlation}

This section addresses the question:

\noindent\textit{What is the cross-asset correlation of daily returns, and which assets offer the strongest diversification potential?}

\label{appendix-a3-sql-question-3}




\begin{Shaded}
\begin{Highlighting}[]
\CommentTok{\# Q3: Pairwise Correlation of Daily Returns {-}{-} SISON, Aaron Joshua Estacio}

\NormalTok{returns\_wide }\OtherTok{\textless{}{-}}\NormalTok{ fin\_data }\SpecialCharTok{|\textgreater{}}
  \FunctionTok{select}\NormalTok{(Date, Asset, Daily\_Return) }\SpecialCharTok{|\textgreater{}}
  \FunctionTok{collect}\NormalTok{() }\SpecialCharTok{|\textgreater{}}
  \FunctionTok{pivot\_wider}\NormalTok{(}\AttributeTok{names\_from =}\NormalTok{ Asset, }\AttributeTok{values\_from =}\NormalTok{ Daily\_Return)}

\NormalTok{cor\_matrix }\OtherTok{\textless{}{-}}\NormalTok{ returns\_wide }\SpecialCharTok{|\textgreater{}}
  \FunctionTok{select}\NormalTok{(AAPL, BDO, BTC, SPY) }\SpecialCharTok{|\textgreater{}}
  \FunctionTok{cor}\NormalTok{(}\AttributeTok{use =} \StringTok{"pairwise.complete.obs"}\NormalTok{)}

\NormalTok{cor\_long }\OtherTok{\textless{}{-}} \FunctionTok{as.data.frame}\NormalTok{(cor\_matrix) }\SpecialCharTok{|\textgreater{}}
\NormalTok{  tibble}\SpecialCharTok{::}\FunctionTok{rownames\_to\_column}\NormalTok{(}\StringTok{"Asset1"}\NormalTok{) }\SpecialCharTok{|\textgreater{}}
\NormalTok{  tidyr}\SpecialCharTok{::}\FunctionTok{pivot\_longer}\NormalTok{(}\SpecialCharTok{{-}}\NormalTok{Asset1, }\AttributeTok{names\_to =} \StringTok{"Asset2"}\NormalTok{, }\AttributeTok{values\_to =} \StringTok{"Correlation"}\NormalTok{) }\SpecialCharTok{|\textgreater{}}
  \FunctionTok{mutate}\NormalTok{(}\AttributeTok{Asset1 =} \FunctionTok{factor}\NormalTok{(Asset1, }\AttributeTok{levels =} \FunctionTok{rev}\NormalTok{(}\FunctionTok{c}\NormalTok{(}\StringTok{"AAPL"}\NormalTok{, }\StringTok{"BDO"}\NormalTok{, }\StringTok{"BTC"}\NormalTok{, }\StringTok{"SPY"}\NormalTok{))),}
         \AttributeTok{Asset2 =} \FunctionTok{factor}\NormalTok{(Asset2, }\AttributeTok{levels =} \FunctionTok{c}\NormalTok{(}\StringTok{"AAPL"}\NormalTok{, }\StringTok{"BDO"}\NormalTok{, }\StringTok{"BTC"}\NormalTok{, }\StringTok{"SPY"}\NormalTok{)))}

\FunctionTok{ggplot}\NormalTok{(cor\_long, }\FunctionTok{aes}\NormalTok{(}\AttributeTok{x =}\NormalTok{ Asset2, }\AttributeTok{y =}\NormalTok{ Asset1, }\AttributeTok{fill =}\NormalTok{ Correlation)) }\SpecialCharTok{+}
  \FunctionTok{geom\_tile}\NormalTok{(}\AttributeTok{color =} \StringTok{"white"}\NormalTok{, }\AttributeTok{linewidth =} \DecValTok{1}\NormalTok{) }\SpecialCharTok{+}
  \FunctionTok{geom\_text}\NormalTok{(}\FunctionTok{aes}\NormalTok{(}\AttributeTok{label =} \FunctionTok{sprintf}\NormalTok{(}\StringTok{"\%.3f"}\NormalTok{, Correlation)), }\AttributeTok{size =} \FloatTok{3.5}\NormalTok{) }\SpecialCharTok{+}
  \FunctionTok{scale\_fill\_gradient2}\NormalTok{(}\AttributeTok{low =} \StringTok{"\#1DC9A4"}\NormalTok{, }\AttributeTok{mid =} \StringTok{"\#F2F2F2"}\NormalTok{, }\AttributeTok{high =} \StringTok{"\#2E45B8"}\NormalTok{,}
                       \AttributeTok{midpoint =} \FloatTok{0.4}\NormalTok{, }\AttributeTok{limits =} \FunctionTok{c}\NormalTok{(}\DecValTok{0}\NormalTok{, }\DecValTok{1}\NormalTok{)) }\SpecialCharTok{+}
  \FunctionTok{labs}\NormalTok{(}\AttributeTok{title =} \StringTok{"Pairwise Correlation of Daily Returns"}\NormalTok{,}
       \AttributeTok{subtitle =} \StringTok{"Pearson correlation coefficients, 2020 to 2025"}\NormalTok{,}
       \AttributeTok{x =} \ConstantTok{NULL}\NormalTok{, }\AttributeTok{y =} \ConstantTok{NULL}\NormalTok{, }\AttributeTok{fill =} \StringTok{"Correlation"}\NormalTok{) }\SpecialCharTok{+}
  \FunctionTok{theme\_minimal}\NormalTok{(}\AttributeTok{base\_size =} \DecValTok{11}\NormalTok{) }\SpecialCharTok{+}
  \FunctionTok{theme}\NormalTok{(}\AttributeTok{panel.grid =} \FunctionTok{element\_blank}\NormalTok{(),}
        \AttributeTok{axis.ticks =} \FunctionTok{element\_blank}\NormalTok{(),}
        \AttributeTok{plot.title =} \FunctionTok{element\_text}\NormalTok{(}\AttributeTok{face =} \StringTok{"bold"}\NormalTok{, }\AttributeTok{size =} \DecValTok{13}\NormalTok{),}
        \AttributeTok{plot.subtitle =} \FunctionTok{element\_text}\NormalTok{(}\AttributeTok{size =} \DecValTok{9}\NormalTok{, }\AttributeTok{color =} \StringTok{"\#595959"}\NormalTok{),}
        \AttributeTok{legend.position =} \StringTok{"right"}\NormalTok{)}
\end{Highlighting}
\end{Shaded}

\begin{table}[H]
\centering
\caption{Q3 Correlation Matrix Output}
\begin{tabular}{lrrrr}
\toprule
      & AAPL   & BDO    & BTC    & SPY \\
\midrule
AAPL & 1.0000 & 0.1044 & 0.2836 & 0.7839 \\
BDO  & 0.1044 & 1.0000 & 0.0507 & 0.1235 \\
BTC  & 0.2836 & 0.0507 & 1.0000 & 0.3825 \\
SPY  & 0.7839 & 0.1235 & 0.3825 & 1.0000 \\
\bottomrule
\end{tabular}
\end{table}

\begin{Shaded}
\begin{Highlighting}[]
\FunctionTok{ggsave}\NormalTok{(}\StringTok{"../reports/figures/fig\_q3.pdf"}\NormalTok{, }\AttributeTok{width =} \DecValTok{6}\NormalTok{, }\AttributeTok{height =} \FloatTok{4.5}\NormalTok{, }\AttributeTok{device =} \StringTok{"pdf"}\NormalTok{)}
\end{Highlighting}
\end{Shaded}

Table~\ref{tab:q3} computes the pairwise Pearson correlation matrix. The intersection of AAPL and SPY produces the highest correlation at 0.784, reflecting Apple's heavy weighting in the index. BDO's correlations with all other assets fall between 0.05 and 0.12, reflecting the structural isolation of the Philippine banking sector. BTC holds a 0.383 correlation with SPY, indicating it partially behaves as a risk asset.

\subsection{Q4. Asset Ranking by Total Return}

This section addresses the question:

\noindent\textit{How do the selected assets rank by total return from 2020 to 2025?}

\label{appendix-a4-sql-question-4}




\begin{Shaded}
\begin{Highlighting}[]
\CommentTok{\# Q4: Asset Ranking by Total Return {-}{-} CRUZ, Ricardo Miguel Inigo}

\NormalTok{q4\_raw }\OtherTok{\textless{}{-}}\NormalTok{ fin\_data }\SpecialCharTok{|\textgreater{}}
  \FunctionTok{select}\NormalTok{(Date, Asset, Close, Daily\_Return) }\SpecialCharTok{|\textgreater{}}
  \FunctionTok{collect}\NormalTok{() }\SpecialCharTok{|\textgreater{}}
  \FunctionTok{mutate}\NormalTok{(}\AttributeTok{Date =} \FunctionTok{as.Date}\NormalTok{(Date))}

\NormalTok{q4 }\OtherTok{\textless{}{-}}\NormalTok{ q4\_raw }\SpecialCharTok{|\textgreater{}}
  \FunctionTok{filter}\NormalTok{(}\SpecialCharTok{!}\FunctionTok{is.na}\NormalTok{(Close), }\SpecialCharTok{!}\FunctionTok{is.na}\NormalTok{(Daily\_Return)) }\SpecialCharTok{|\textgreater{}}
  \FunctionTok{group\_by}\NormalTok{(Asset) }\SpecialCharTok{|\textgreater{}}
  \FunctionTok{arrange}\NormalTok{(Date) }\SpecialCharTok{|\textgreater{}}
  \FunctionTok{summarise}\NormalTok{(}
    \AttributeTok{Start\_Date =} \FunctionTok{min}\NormalTok{(Date),}
    \AttributeTok{End\_Date =} \FunctionTok{max}\NormalTok{(Date),}
    \AttributeTok{Start\_Close =} \FunctionTok{first}\NormalTok{(Close),}
    \AttributeTok{End\_Close =} \FunctionTok{last}\NormalTok{(Close),}
    \AttributeTok{Total\_Return =}\NormalTok{ (End\_Close }\SpecialCharTok{/}\NormalTok{ Start\_Close) }\SpecialCharTok{{-}} \DecValTok{1}\NormalTok{,}
    \AttributeTok{Annualized\_Return =}\NormalTok{ (}\DecValTok{1} \SpecialCharTok{+}\NormalTok{ Total\_Return)}\SpecialCharTok{\^{}}\NormalTok{(}\FloatTok{365.25} \SpecialCharTok{/} \FunctionTok{as.numeric}\NormalTok{(End\_Date }\SpecialCharTok{{-}}\NormalTok{ Start\_Date)) }\SpecialCharTok{{-}} \DecValTok{1}\NormalTok{,}
    \AttributeTok{Mean\_Daily\_Return =} \FunctionTok{mean}\NormalTok{(Daily\_Return, }\AttributeTok{na.rm =} \ConstantTok{TRUE}\NormalTok{),}
    \AttributeTok{SD\_Daily\_Return =} \FunctionTok{sd}\NormalTok{(Daily\_Return, }\AttributeTok{na.rm =} \ConstantTok{TRUE}\NormalTok{),}
    \AttributeTok{Observations =} \FunctionTok{n}\NormalTok{(),}
    \AttributeTok{.groups =} \StringTok{"drop"}
\NormalTok{  ) }\SpecialCharTok{|\textgreater{}}
  \FunctionTok{arrange}\NormalTok{(}\FunctionTok{desc}\NormalTok{(Total\_Return))}

\NormalTok{q4\_plot }\OtherTok{\textless{}{-}}\NormalTok{ q4 }\SpecialCharTok{|\textgreater{}}
  \FunctionTok{mutate}\NormalTok{(}\AttributeTok{Asset =} \FunctionTok{factor}\NormalTok{(Asset, }\AttributeTok{levels =}\NormalTok{ q4}\SpecialCharTok{$}\NormalTok{Asset[}\FunctionTok{order}\NormalTok{(q4}\SpecialCharTok{$}\NormalTok{Total\_Return)]))}

\FunctionTok{ggplot}\NormalTok{(q4\_plot, }\FunctionTok{aes}\NormalTok{(}\AttributeTok{x =}\NormalTok{ Total\_Return, }\AttributeTok{y =}\NormalTok{ Asset, }\AttributeTok{fill =}\NormalTok{ Asset)) }\SpecialCharTok{+}
  \FunctionTok{geom\_col}\NormalTok{(}\AttributeTok{width =} \FloatTok{0.6}\NormalTok{, }\AttributeTok{show.legend =} \ConstantTok{FALSE}\NormalTok{) }\SpecialCharTok{+}
  \FunctionTok{geom\_vline}\NormalTok{(}\AttributeTok{xintercept =} \DecValTok{0}\NormalTok{, }\AttributeTok{linewidth =} \FloatTok{0.3}\NormalTok{, }\AttributeTok{color =} \StringTok{"\#595959"}\NormalTok{) }\SpecialCharTok{+}
  \FunctionTok{geom\_text}\NormalTok{(}\FunctionTok{aes}\NormalTok{(}\AttributeTok{x =}\NormalTok{ Total\_Return }\SpecialCharTok{/} \DecValTok{2}\NormalTok{,}
                \AttributeTok{label =} \FunctionTok{sprintf}\NormalTok{(}\StringTok{"\%+.1f\%\%"}\NormalTok{, Total\_Return }\SpecialCharTok{*} \DecValTok{100}\NormalTok{)),}
            \AttributeTok{hjust =} \FloatTok{0.5}\NormalTok{, }\AttributeTok{size =} \FloatTok{3.2}\NormalTok{, }\AttributeTok{color =} \StringTok{"white"}\NormalTok{, }\AttributeTok{fontface =} \StringTok{"bold"}\NormalTok{) }\SpecialCharTok{+}
  \FunctionTok{scale\_fill\_manual}\NormalTok{(}\AttributeTok{values =} \FunctionTok{c}\NormalTok{(}\AttributeTok{AAPL =} \StringTok{"\#555555"}\NormalTok{, }\AttributeTok{BDO =} \StringTok{"\#003D6B"}\NormalTok{,}
                                \AttributeTok{BTC =} \StringTok{"\#F7931A"}\NormalTok{, }\AttributeTok{SPY =} \StringTok{"\#00733E"}\NormalTok{)) }\SpecialCharTok{+}
  \FunctionTok{scale\_x\_continuous}\NormalTok{(}\AttributeTok{labels =}\NormalTok{ scales}\SpecialCharTok{::}\FunctionTok{percent\_format}\NormalTok{(),}
                     \AttributeTok{expand =} \FunctionTok{expansion}\NormalTok{(}\AttributeTok{mult =} \FunctionTok{c}\NormalTok{(}\FloatTok{0.15}\NormalTok{, }\FloatTok{0.15}\NormalTok{))) }\SpecialCharTok{+}
  \FunctionTok{labs}\NormalTok{(}\AttributeTok{title =} \StringTok{"Asset Ranking by Total Return"}\NormalTok{,}
       \AttributeTok{subtitle =} \StringTok{"Total price return from 2020 to 2025"}\NormalTok{,}
       \AttributeTok{x =} \StringTok{"Total Return"}\NormalTok{, }\AttributeTok{y =} \ConstantTok{NULL}\NormalTok{) }\SpecialCharTok{+}
  \FunctionTok{theme\_minimal}\NormalTok{(}\AttributeTok{base\_size =} \DecValTok{11}\NormalTok{) }\SpecialCharTok{+}
  \FunctionTok{theme}\NormalTok{(}\AttributeTok{panel.grid.major.y =} \FunctionTok{element\_blank}\NormalTok{(),}
        \AttributeTok{panel.grid.minor =} \FunctionTok{element\_blank}\NormalTok{(),}
        \AttributeTok{axis.ticks =} \FunctionTok{element\_blank}\NormalTok{(),}
        \AttributeTok{plot.title =} \FunctionTok{element\_text}\NormalTok{(}\AttributeTok{face =} \StringTok{"bold"}\NormalTok{, }\AttributeTok{size =} \DecValTok{13}\NormalTok{),}
        \AttributeTok{plot.subtitle =} \FunctionTok{element\_text}\NormalTok{(}\AttributeTok{size =} \DecValTok{9}\NormalTok{, }\AttributeTok{color =} \StringTok{"\#595959"}\NormalTok{))}
\end{Highlighting}
\end{Shaded}

\begin{table}[H]
\centering
\caption{Q4 Data Output}
\begin{tabular}{lrrrrrrrrr}
\toprule
Asset & Start\_Date & End\_Date & Start\_Close & End\_Close & Total\_Return & Annualized\_Return & Mean\_Daily\_Return & SD\_Daily\_Return & Obs \\
\midrule
BTC  & 2020-01-02 & 2025-12-31 & 6985.47 & 87508.83 & 11.5273 & 0.5244 & 0.0017 & 0.0319 & 2191 \\
AAPL & 2020-01-03 & 2025-12-31 & 71.63  & 271.36  & 2.7884 & 0.2489 & 0.0011 & 0.0200 & 1507 \\
SPY  & 2020-01-03 & 2025-12-31 & 293.88 & 678.32  & 1.3082 & 0.1498 & 0.0006 & 0.0131 & 1507 \\
BDO  & 2020-01-03 & 2025-12-29 & 128.17 & 134.60  & 0.0502 & 0.0082 & 0.0003 & 0.0226 & 1467 \\
\bottomrule
\end{tabular}
\end{table}

\begin{Shaded}
\begin{Highlighting}[]
\FunctionTok{ggsave}\NormalTok{(}\StringTok{"../reports/figures/fig\_q4.pdf"}\NormalTok{, }\AttributeTok{width =} \DecValTok{7}\NormalTok{, }\AttributeTok{height =} \FloatTok{3.5}\NormalTok{, }\AttributeTok{device =} \StringTok{"pdf"}\NormalTok{)}
\end{Highlighting}
\end{Shaded}

Table~\ref{tab:q4} ranks the assets by total absolute return from 2020 to 2025. Bitcoin delivered a 1,153\% total return over the six-year period. Apple returned 279\%, the S\&P 500 returned 131\%, and BDO Unibank stagnated at a 5.0\% total return. The query isolates the magnitude of outperformance during the zero-interest-rate policy era.

\subsection{Q5. Best and Worst Trading Days}

This section addresses the question:

\noindent\textit{What were the best and worst single-day returns for each asset, and what do they reveal about structural tail risk?}

\label{appendix-a5-sql-question-5}

\emph{Pertains to Section 6: SQL-Based Financial Investigation}

\textbf{Objective:} To identify the extreme single-day gains and losses
for each asset. \textbf{Purpose of Code:} Extracts the maximum and
minimum daily returns for each instrument to assess tail risk.
\textbf{Expected Output:} A grouped bar chart displaying the best and
worst trading days.

\begin{Shaded}
\begin{Highlighting}[]
\CommentTok{\# Q5: Best and Worst Trading Days {-}{-} CRUZ, Ricardo Miguel Inigo}

\NormalTok{q5\_raw }\OtherTok{\textless{}{-}}\NormalTok{ fin\_data }\SpecialCharTok{|\textgreater{}}
  \FunctionTok{select}\NormalTok{(Date, Asset, Daily\_Return) }\SpecialCharTok{|\textgreater{}}
  \FunctionTok{collect}\NormalTok{() }\SpecialCharTok{|\textgreater{}}
  \FunctionTok{mutate}\NormalTok{(}\AttributeTok{Date =} \FunctionTok{as.Date}\NormalTok{(Date))}

\NormalTok{q5 }\OtherTok{\textless{}{-}}\NormalTok{ q5\_raw }\SpecialCharTok{|\textgreater{}}
  \FunctionTok{filter}\NormalTok{(}\SpecialCharTok{!}\FunctionTok{is.na}\NormalTok{(Daily\_Return)) }\SpecialCharTok{|\textgreater{}}
  \FunctionTok{group\_by}\NormalTok{(Asset) }\SpecialCharTok{|\textgreater{}}
  \FunctionTok{summarise}\NormalTok{(}
    \AttributeTok{Best\_Date =}\NormalTok{ Date[}\FunctionTok{which.max}\NormalTok{(Daily\_Return)],}
    \AttributeTok{Best\_Daily\_Return =} \FunctionTok{max}\NormalTok{(Daily\_Return, }\AttributeTok{na.rm =} \ConstantTok{TRUE}\NormalTok{),}
    \AttributeTok{Worst\_Date =}\NormalTok{ Date[}\FunctionTok{which.min}\NormalTok{(Daily\_Return)],}
    \AttributeTok{Worst\_Daily\_Return =} \FunctionTok{min}\NormalTok{(Daily\_Return, }\AttributeTok{na.rm =} \ConstantTok{TRUE}\NormalTok{),}
    \AttributeTok{Return\_Range =}\NormalTok{ Best\_Daily\_Return }\SpecialCharTok{{-}}\NormalTok{ Worst\_Daily\_Return,}
    \AttributeTok{Observations =} \FunctionTok{n}\NormalTok{(),}
    \AttributeTok{.groups =} \StringTok{"drop"}
\NormalTok{  ) }\SpecialCharTok{|\textgreater{}}
  \FunctionTok{arrange}\NormalTok{(Worst\_Daily\_Return)}

\NormalTok{q5\_plot }\OtherTok{\textless{}{-}}\NormalTok{ q5 }\SpecialCharTok{|\textgreater{}}
  \FunctionTok{select}\NormalTok{(Asset, Best\_Daily\_Return, Worst\_Daily\_Return) }\SpecialCharTok{|\textgreater{}}
  \FunctionTok{pivot\_longer}\NormalTok{(}\AttributeTok{cols =} \FunctionTok{c}\NormalTok{(Best\_Daily\_Return, Worst\_Daily\_Return),}
               \AttributeTok{names\_to =} \StringTok{"Trading\_Day\_Type"}\NormalTok{,}
               \AttributeTok{values\_to =} \StringTok{"Daily\_Return"}\NormalTok{) }\SpecialCharTok{|\textgreater{}}
  \FunctionTok{mutate}\NormalTok{(}
    \AttributeTok{Trading\_Day\_Type =} \FunctionTok{recode}\NormalTok{(Trading\_Day\_Type,}
                              \AttributeTok{Best\_Daily\_Return =} \StringTok{"Best Trading Day"}\NormalTok{,}
                              \AttributeTok{Worst\_Daily\_Return =} \StringTok{"Worst Trading Day"}\NormalTok{)}
\NormalTok{  )}

\FunctionTok{ggplot}\NormalTok{(q5\_plot, }\FunctionTok{aes}\NormalTok{(}\AttributeTok{x =}\NormalTok{ Asset, }\AttributeTok{y =}\NormalTok{ Daily\_Return, }\AttributeTok{fill =}\NormalTok{ Trading\_Day\_Type)) }\SpecialCharTok{+}
  \FunctionTok{geom\_col}\NormalTok{(}\AttributeTok{position =} \FunctionTok{position\_dodge}\NormalTok{(}\AttributeTok{width =} \FloatTok{0.7}\NormalTok{), }\AttributeTok{width =} \FloatTok{0.6}\NormalTok{) }\SpecialCharTok{+}
  \FunctionTok{geom\_hline}\NormalTok{(}\AttributeTok{yintercept =} \DecValTok{0}\NormalTok{, }\AttributeTok{linewidth =} \FloatTok{0.3}\NormalTok{, }\AttributeTok{color =} \StringTok{"\#595959"}\NormalTok{) }\SpecialCharTok{+}
  \FunctionTok{scale\_fill\_manual}\NormalTok{(}\AttributeTok{values =} \FunctionTok{c}\NormalTok{(}\StringTok{"Best Trading Day"} \OtherTok{=} \StringTok{"\#2E45B8"}\NormalTok{,}
                                \StringTok{"Worst Trading Day"} \OtherTok{=} \StringTok{"\#C91D42"}\NormalTok{)) }\SpecialCharTok{+}
  \FunctionTok{scale\_y\_continuous}\NormalTok{(}\AttributeTok{labels =}\NormalTok{ scales}\SpecialCharTok{::}\FunctionTok{percent\_format}\NormalTok{()) }\SpecialCharTok{+}
  \FunctionTok{labs}\NormalTok{(}\AttributeTok{title =} \StringTok{"Best and Worst Single{-}Day Returns"}\NormalTok{,}
       \AttributeTok{subtitle =} \StringTok{"Largest one{-}day gains and losses by asset, 2020 to 2025"}\NormalTok{,}
       \AttributeTok{x =} \ConstantTok{NULL}\NormalTok{, }\AttributeTok{y =} \StringTok{"Daily Return"}\NormalTok{, }\AttributeTok{fill =} \ConstantTok{NULL}\NormalTok{) }\SpecialCharTok{+}
  \FunctionTok{theme\_minimal}\NormalTok{(}\AttributeTok{base\_size =} \DecValTok{11}\NormalTok{) }\SpecialCharTok{+}
  \FunctionTok{theme}\NormalTok{(}\AttributeTok{panel.grid.major.x =} \FunctionTok{element\_blank}\NormalTok{(),}
        \AttributeTok{panel.grid.minor =} \FunctionTok{element\_blank}\NormalTok{(),}
        \AttributeTok{axis.ticks =} \FunctionTok{element\_blank}\NormalTok{(),}
        \AttributeTok{plot.title =} \FunctionTok{element\_text}\NormalTok{(}\AttributeTok{face =} \StringTok{"bold"}\NormalTok{, }\AttributeTok{size =} \DecValTok{13}\NormalTok{),}
        \AttributeTok{plot.subtitle =} \FunctionTok{element\_text}\NormalTok{(}\AttributeTok{size =} \DecValTok{9}\NormalTok{, }\AttributeTok{color =} \StringTok{"\#595959"}\NormalTok{),}
        \AttributeTok{legend.position =} \StringTok{"bottom"}\NormalTok{)}
\end{Highlighting}
\end{Shaded}



\begin{Shaded}
\begin{Highlighting}[]
\FunctionTok{ggsave}\NormalTok{(}\StringTok{"../reports/figures/fig\_q5.pdf"}\NormalTok{, }\AttributeTok{width =} \DecValTok{7}\NormalTok{, }\AttributeTok{height =} \DecValTok{4}\NormalTok{, }\AttributeTok{device =} \StringTok{"pdf"}\NormalTok{)}
\end{Highlighting}
\end{Shaded}

Table~\ref{tab:q5} isolates the structural tail risk by extracting the best and worst single-day returns. BTC recorded the widest spread between its best trading day of +18.7\% and its worst of -37.2\%. SPY exhibited the narrowest range (+10.5\% to -10.9\%). AAPL and BDO recorded maximum drawdowns of -12.9\% and -22.7\% respectively.

\subsection{Q6. Moving Average Crossover Analysis}

This section addresses the question:

\noindent\textit{How frequently do short-term price trends deviate from medium-term baselines across different assets?}

\label{appendix-a6-sql-question-6}

\emph{Pertains to Section 6: SQL-Based Financial Investigation}

\textbf{Objective:} To analyze momentum regime shifts using moving
averages. \textbf{Purpose of Code:} Implements a 20-day and 60-day
simple moving average crossover strategy. \textbf{Expected Output:} A
multi-panel line chart showing crossover events and momentum trends.

\begin{Shaded}
\begin{Highlighting}[]
\CommentTok{\# Q6: Moving Average Crossover Analysis {-}{-} CRUZ, Ricardo Miguel Inigo}

\NormalTok{q6\_raw }\OtherTok{\textless{}{-}}\NormalTok{ fin\_data }\SpecialCharTok{|\textgreater{}}
  \FunctionTok{select}\NormalTok{(Date, Asset, Close, Daily\_Return) }\SpecialCharTok{|\textgreater{}}
  \FunctionTok{collect}\NormalTok{() }\SpecialCharTok{|\textgreater{}}
  \FunctionTok{mutate}\NormalTok{(}\AttributeTok{Date =} \FunctionTok{as.Date}\NormalTok{(Date))}

\NormalTok{q6\_ma }\OtherTok{\textless{}{-}}\NormalTok{ q6\_raw }\SpecialCharTok{|\textgreater{}}
  \FunctionTok{filter}\NormalTok{(}\SpecialCharTok{!}\FunctionTok{is.na}\NormalTok{(Close)) }\SpecialCharTok{|\textgreater{}}
  \FunctionTok{group\_by}\NormalTok{(Asset) }\SpecialCharTok{|\textgreater{}}
  \FunctionTok{arrange}\NormalTok{(Date) }\SpecialCharTok{|\textgreater{}}
  \FunctionTok{mutate}\NormalTok{(}
    \AttributeTok{MA\_20 =} \FunctionTok{as.numeric}\NormalTok{(stats}\SpecialCharTok{::}\FunctionTok{filter}\NormalTok{(Close, }\FunctionTok{rep}\NormalTok{(}\DecValTok{1} \SpecialCharTok{/} \DecValTok{20}\NormalTok{, }\DecValTok{20}\NormalTok{), }\AttributeTok{sides =} \DecValTok{1}\NormalTok{)),}
    \AttributeTok{MA\_60 =} \FunctionTok{as.numeric}\NormalTok{(stats}\SpecialCharTok{::}\FunctionTok{filter}\NormalTok{(Close, }\FunctionTok{rep}\NormalTok{(}\DecValTok{1} \SpecialCharTok{/} \DecValTok{60}\NormalTok{, }\DecValTok{60}\NormalTok{), }\AttributeTok{sides =} \DecValTok{1}\NormalTok{)),}
    \AttributeTok{Signal =} \FunctionTok{case\_when}\NormalTok{(}
\NormalTok{      MA\_20 }\SpecialCharTok{\textgreater{}}\NormalTok{ MA\_60 }\SpecialCharTok{\textasciitilde{}} \StringTok{"Bullish"}\NormalTok{,}
\NormalTok{      MA\_20 }\SpecialCharTok{\textless{}}\NormalTok{ MA\_60 }\SpecialCharTok{\textasciitilde{}} \StringTok{"Bearish"}\NormalTok{,}
      \ConstantTok{TRUE} \SpecialCharTok{\textasciitilde{}} \StringTok{"Neutral"}
\NormalTok{    )}
\NormalTok{  ) }\SpecialCharTok{|\textgreater{}}
  \FunctionTok{ungroup}\NormalTok{()}

\NormalTok{q6 }\OtherTok{\textless{}{-}}\NormalTok{ q6\_ma }\SpecialCharTok{|\textgreater{}}
  \FunctionTok{filter}\NormalTok{(}\SpecialCharTok{!}\FunctionTok{is.na}\NormalTok{(MA\_20), }\SpecialCharTok{!}\FunctionTok{is.na}\NormalTok{(MA\_60), }\SpecialCharTok{!}\FunctionTok{is.na}\NormalTok{(Daily\_Return)) }\SpecialCharTok{|\textgreater{}}
  \FunctionTok{group\_by}\NormalTok{(Asset, Signal) }\SpecialCharTok{|\textgreater{}}
  \FunctionTok{summarise}\NormalTok{(}
    \AttributeTok{Mean\_Daily\_Return =} \FunctionTok{mean}\NormalTok{(Daily\_Return, }\AttributeTok{na.rm =} \ConstantTok{TRUE}\NormalTok{),}
    \AttributeTok{SD\_Daily\_Return =} \FunctionTok{sd}\NormalTok{(Daily\_Return, }\AttributeTok{na.rm =} \ConstantTok{TRUE}\NormalTok{),}
    \AttributeTok{Observations =} \FunctionTok{n}\NormalTok{(),}
    \AttributeTok{.groups =} \StringTok{"drop"}
\NormalTok{  ) }\SpecialCharTok{|\textgreater{}}
  \FunctionTok{arrange}\NormalTok{(Asset, }\FunctionTok{desc}\NormalTok{(Mean\_Daily\_Return))}

\NormalTok{q6\_latest\_signal }\OtherTok{\textless{}{-}}\NormalTok{ q6\_ma }\SpecialCharTok{|\textgreater{}}
  \FunctionTok{filter}\NormalTok{(}\SpecialCharTok{!}\FunctionTok{is.na}\NormalTok{(MA\_20), }\SpecialCharTok{!}\FunctionTok{is.na}\NormalTok{(MA\_60)) }\SpecialCharTok{|\textgreater{}}
  \FunctionTok{group\_by}\NormalTok{(Asset) }\SpecialCharTok{|\textgreater{}}
  \FunctionTok{slice\_max}\NormalTok{(Date, }\AttributeTok{n =} \DecValTok{1}\NormalTok{, }\AttributeTok{with\_ties =} \ConstantTok{FALSE}\NormalTok{) }\SpecialCharTok{|\textgreater{}}
  \FunctionTok{ungroup}\NormalTok{() }\SpecialCharTok{|\textgreater{}}
  \FunctionTok{select}\NormalTok{(Asset, Date, Close, MA\_20, MA\_60, Signal) }\SpecialCharTok{|\textgreater{}}
  \FunctionTok{arrange}\NormalTok{(Asset)}

\NormalTok{q6\_plot }\OtherTok{\textless{}{-}}\NormalTok{ q6\_ma }\SpecialCharTok{|\textgreater{}}
  \FunctionTok{filter}\NormalTok{(Date }\SpecialCharTok{\textgreater{}=} \FunctionTok{as.Date}\NormalTok{(}\StringTok{"2024{-}01{-}01"}\NormalTok{))}

\FunctionTok{ggplot}\NormalTok{(q6\_plot, }\FunctionTok{aes}\NormalTok{(}\AttributeTok{x =}\NormalTok{ Date)) }\SpecialCharTok{+}
  \FunctionTok{geom\_line}\NormalTok{(}\FunctionTok{aes}\NormalTok{(}\AttributeTok{y =}\NormalTok{ Close), }\AttributeTok{linewidth =} \FloatTok{0.4}\NormalTok{, }\AttributeTok{color =} \StringTok{"\#595959"}\NormalTok{) }\SpecialCharTok{+}
  \FunctionTok{geom\_line}\NormalTok{(}\FunctionTok{aes}\NormalTok{(}\AttributeTok{y =}\NormalTok{ MA\_20), }\AttributeTok{linewidth =} \FloatTok{0.5}\NormalTok{, }\AttributeTok{color =} \StringTok{"\#2E45B8"}\NormalTok{) }\SpecialCharTok{+}
  \FunctionTok{geom\_line}\NormalTok{(}\FunctionTok{aes}\NormalTok{(}\AttributeTok{y =}\NormalTok{ MA\_60), }\AttributeTok{linewidth =} \FloatTok{0.5}\NormalTok{, }\AttributeTok{color =} \StringTok{"\#C91D42"}\NormalTok{) }\SpecialCharTok{+}
  \FunctionTok{facet\_wrap}\NormalTok{(}\SpecialCharTok{\textasciitilde{}}\NormalTok{ Asset, }\AttributeTok{scales =} \StringTok{"free\_y"}\NormalTok{) }\SpecialCharTok{+}
  \FunctionTok{labs}\NormalTok{(}\AttributeTok{title =} \StringTok{"20{-}Day and 60{-}Day Moving Average Crossover"}\NormalTok{,}
       \AttributeTok{subtitle =} \StringTok{"Bullish signal when 20{-}day MA is above 60{-}day MA, 2024 to 2025"}\NormalTok{,}
       \AttributeTok{x =} \ConstantTok{NULL}\NormalTok{, }\AttributeTok{y =} \StringTok{"Closing Price"}\NormalTok{) }\SpecialCharTok{+}
  \FunctionTok{theme\_minimal}\NormalTok{(}\AttributeTok{base\_size =} \DecValTok{11}\NormalTok{) }\SpecialCharTok{+}
  \FunctionTok{theme}\NormalTok{(}\AttributeTok{panel.grid.minor =} \FunctionTok{element\_blank}\NormalTok{(),}
        \AttributeTok{axis.ticks =} \FunctionTok{element\_blank}\NormalTok{(),}
        \AttributeTok{plot.title =} \FunctionTok{element\_text}\NormalTok{(}\AttributeTok{face =} \StringTok{"bold"}\NormalTok{, }\AttributeTok{size =} \DecValTok{13}\NormalTok{),}
        \AttributeTok{plot.subtitle =} \FunctionTok{element\_text}\NormalTok{(}\AttributeTok{size =} \DecValTok{9}\NormalTok{, }\AttributeTok{color =} \StringTok{"\#595959"}\NormalTok{))}
\end{Highlighting}
\end{Shaded}



\begin{Shaded}
\begin{Highlighting}[]
\FunctionTok{ggsave}\NormalTok{(}\StringTok{"../reports/figures/fig\_q6.pdf"}\NormalTok{, }\AttributeTok{width =} \DecValTok{7}\NormalTok{, }\AttributeTok{height =} \FloatTok{4.5}\NormalTok{, }\AttributeTok{device =} \StringTok{"pdf"}\NormalTok{)}
\end{Highlighting}
\end{Shaded}

Table~\ref{tab:q6} summarizes the mean returns conditional on moving average crossover signals. AAPL and BDO recorded exactly 6 crossover events, SPY recorded 8 intersections, and BTC generated 15 distinct crossovers. The query quantifies the frequency of momentum regime switching across the asset sample.

\subsection{Q7. Volume and Volatility Relationship}

This section addresses the question:

\noindent\textit{Does elevated trading volume systematically coincide with higher daily price volatility?}

\label{appendix-a7-sql-question-7}

\emph{Pertains to Section 6: SQL-Based Financial Investigation}

\textbf{Objective:} To test whether elevated trading volume
systematically coincides with higher volatility. \textbf{Purpose of
Code:} Partitions trading days into high-volume and normal-volume
regimes and compares their mean absolute returns. \textbf{Expected
Output:} A grouped bar chart comparing volatility across volume regimes.

\begin{Shaded}
\begin{Highlighting}[]
\CommentTok{\# Q7: Volume and Volatility Relationship {-}{-} CRUZ, Ricardo Miguel Inigo}

\NormalTok{volume\_lookup }\OtherTok{\textless{}{-}}\NormalTok{ assets\_all }\SpecialCharTok{|\textgreater{}}
  \FunctionTok{select}\NormalTok{(Date, Asset, Volume) }\SpecialCharTok{|\textgreater{}}
  \FunctionTok{mutate}\NormalTok{(}
    \AttributeTok{Date =} \FunctionTok{as.Date}\NormalTok{(Date),}
    \AttributeTok{Volume =} \FunctionTok{as.numeric}\NormalTok{(Volume)}
\NormalTok{  )}

\NormalTok{q7\_raw }\OtherTok{\textless{}{-}}\NormalTok{ fin\_data }\SpecialCharTok{|\textgreater{}}
  \FunctionTok{select}\NormalTok{(Date, Asset, Daily\_Return) }\SpecialCharTok{|\textgreater{}}
  \FunctionTok{collect}\NormalTok{() }\SpecialCharTok{|\textgreater{}}
  \FunctionTok{mutate}\NormalTok{(}
    \AttributeTok{Date =} \FunctionTok{as.Date}\NormalTok{(Date),}
    \AttributeTok{Daily\_Return =} \FunctionTok{as.numeric}\NormalTok{(Daily\_Return)}
\NormalTok{  ) }\SpecialCharTok{|\textgreater{}}
  \FunctionTok{left\_join}\NormalTok{(volume\_lookup, }\AttributeTok{by =} \FunctionTok{c}\NormalTok{(}\StringTok{"Date"}\NormalTok{, }\StringTok{"Asset"}\NormalTok{)) }\SpecialCharTok{|\textgreater{}}
  \FunctionTok{filter}\NormalTok{(}
    \SpecialCharTok{!}\FunctionTok{is.na}\NormalTok{(Daily\_Return),}
    \SpecialCharTok{!}\FunctionTok{is.na}\NormalTok{(Volume),}
\NormalTok{    Volume }\SpecialCharTok{\textgreater{}} \DecValTok{0}
\NormalTok{  )}

\NormalTok{q7 }\OtherTok{\textless{}{-}}\NormalTok{ q7\_raw }\SpecialCharTok{|\textgreater{}}
  \FunctionTok{group\_by}\NormalTok{(Asset) }\SpecialCharTok{|\textgreater{}}
  \FunctionTok{mutate}\NormalTok{(}
    \AttributeTok{Volume\_Threshold =} \FunctionTok{quantile}\NormalTok{(Volume, }\FloatTok{0.75}\NormalTok{, }\AttributeTok{na.rm =} \ConstantTok{TRUE}\NormalTok{),}
    \AttributeTok{Volume\_Regime =} \FunctionTok{if\_else}\NormalTok{(}
\NormalTok{      Volume }\SpecialCharTok{\textgreater{}=}\NormalTok{ Volume\_Threshold,}
      \StringTok{"High Volume Days"}\NormalTok{,}
      \StringTok{"Normal Volume Days"}
\NormalTok{    ),}
    \AttributeTok{Absolute\_Return =} \FunctionTok{abs}\NormalTok{(Daily\_Return)}
\NormalTok{  ) }\SpecialCharTok{|\textgreater{}}
  \FunctionTok{ungroup}\NormalTok{() }\SpecialCharTok{|\textgreater{}}
  \FunctionTok{group\_by}\NormalTok{(Asset, Volume\_Regime) }\SpecialCharTok{|\textgreater{}}
  \FunctionTok{summarise}\NormalTok{(}
    \AttributeTok{Observations =} \FunctionTok{n}\NormalTok{(),}
    \AttributeTok{Mean\_Daily\_Return =} \FunctionTok{mean}\NormalTok{(Daily\_Return, }\AttributeTok{na.rm =} \ConstantTok{TRUE}\NormalTok{),}
    \AttributeTok{Mean\_Absolute\_Return =} \FunctionTok{mean}\NormalTok{(Absolute\_Return, }\AttributeTok{na.rm =} \ConstantTok{TRUE}\NormalTok{),}
    \AttributeTok{Volatility =} \FunctionTok{sd}\NormalTok{(Daily\_Return, }\AttributeTok{na.rm =} \ConstantTok{TRUE}\NormalTok{),}
    \AttributeTok{Average\_Volume =} \FunctionTok{mean}\NormalTok{(Volume, }\AttributeTok{na.rm =} \ConstantTok{TRUE}\NormalTok{),}
    \AttributeTok{.groups =} \StringTok{"drop"}
\NormalTok{  ) }\SpecialCharTok{|\textgreater{}}
  \FunctionTok{arrange}\NormalTok{(Asset, Volume\_Regime)}

\NormalTok{q7}
\end{Highlighting}
\end{Shaded}

\begin{verbatim}
# A tibble: 8 x 7
  Asset Volume_Regime      Observations Mean_Daily_Return
  <chr> <chr>                     <int>             <dbl>
1 AAPL  High Volume Days            377         0.00155  
2 AAPL  Normal Volume Days         1130         0.000921 
3 BDO   High Volume Days            367        -0.0000981
4 BDO   Normal Volume Days         1100         0.000426 
5 BTC   High Volume Days            548         0.00168  
6 BTC   Normal Volume Days         1643         0.00165  
7 SPY   High Volume Days            377        -0.00292  
8 SPY   Normal Volume Days         1130         0.00182  
# i 3 more variables: Mean_Absolute_Return <dbl>, Volatility <dbl>,
#   Average_Volume <dbl>
\end{verbatim}

\begin{Shaded}
\begin{Highlighting}[]
\FunctionTok{ggplot}\NormalTok{(q7, }\FunctionTok{aes}\NormalTok{(}\AttributeTok{x =}\NormalTok{ Asset, }\AttributeTok{y =}\NormalTok{ Mean\_Absolute\_Return, }\AttributeTok{fill =}\NormalTok{ Volume\_Regime)) }\SpecialCharTok{+}
  \FunctionTok{geom\_col}\NormalTok{(}\AttributeTok{position =} \FunctionTok{position\_dodge}\NormalTok{(}\AttributeTok{width =} \FloatTok{0.7}\NormalTok{), }\AttributeTok{width =} \FloatTok{0.6}\NormalTok{) }\SpecialCharTok{+}
  \FunctionTok{scale\_y\_continuous}\NormalTok{(}\AttributeTok{labels =}\NormalTok{ scales}\SpecialCharTok{::}\FunctionTok{percent\_format}\NormalTok{()) }\SpecialCharTok{+}
  \FunctionTok{labs}\NormalTok{(}
    \AttributeTok{title =} \StringTok{"Volume and Volatility Relationship"}\NormalTok{,}
    \AttributeTok{subtitle =} \StringTok{"Mean absolute daily return during high{-}volume versus normal{-}volume days, 2020 to 2025"}\NormalTok{,}
    \AttributeTok{x =} \ConstantTok{NULL}\NormalTok{,}
    \AttributeTok{y =} \StringTok{"Mean Absolute Daily Return"}\NormalTok{,}
    \AttributeTok{fill =} \ConstantTok{NULL}
\NormalTok{  ) }\SpecialCharTok{+}
  \FunctionTok{theme\_minimal}\NormalTok{(}\AttributeTok{base\_size =} \DecValTok{11}\NormalTok{) }\SpecialCharTok{+}
  \FunctionTok{theme}\NormalTok{(}
    \AttributeTok{panel.grid.major.x =} \FunctionTok{element\_blank}\NormalTok{(),}
    \AttributeTok{panel.grid.minor =} \FunctionTok{element\_blank}\NormalTok{(),}
    \AttributeTok{axis.ticks =} \FunctionTok{element\_blank}\NormalTok{(),}
    \AttributeTok{plot.title =} \FunctionTok{element\_text}\NormalTok{(}\AttributeTok{face =} \StringTok{"bold"}\NormalTok{, }\AttributeTok{size =} \DecValTok{13}\NormalTok{),}
    \AttributeTok{plot.subtitle =} \FunctionTok{element\_text}\NormalTok{(}\AttributeTok{size =} \DecValTok{9}\NormalTok{, }\AttributeTok{color =} \StringTok{"\#595959"}\NormalTok{),}
    \AttributeTok{legend.position =} \StringTok{"bottom"}
\NormalTok{  )}
\end{Highlighting}
\end{Shaded}



\begin{Shaded}
\begin{Highlighting}[]
\FunctionTok{dir.create}\NormalTok{(}\StringTok{"../reports/figures"}\NormalTok{, }\AttributeTok{recursive =} \ConstantTok{TRUE}\NormalTok{, }\AttributeTok{showWarnings =} \ConstantTok{FALSE}\NormalTok{)}

\FunctionTok{ggsave}\NormalTok{(}\StringTok{"../reports/figures/fig\_q7.pdf"}\NormalTok{, }\AttributeTok{width =} \DecValTok{7}\NormalTok{, }\AttributeTok{height =} \DecValTok{4}\NormalTok{, }\AttributeTok{device =} \StringTok{"pdf"}\NormalTok{)}
\end{Highlighting}
\end{Shaded}

Table~\ref{tab:q7} partitions the descriptive statistics by volume regime. BTC exhibited a mean absolute return of 3.05\% on high-volume days compared to 1.83\% on normal days. SPY showed a proportional expansion from 0.61\% to 1.60\%. The data confirm that elevated trading volume systematically coincides with higher daily price volatility across all asset classes.

\subsection{Q8. VIX Regime Analysis}

This section addresses the question:

\noindent\textit{How does asset behaviour change during high-fear (VIX > 30) versus low-fear (VIX < 20) market regimes?}

\label{appendix-a8-sql-question-8}




\begin{Shaded}
\begin{Highlighting}[]
\CommentTok{\# Q8: VIX Regime Analysis {-}{-} SISON, Aaron Joshua Estacio}

\NormalTok{q8\_raw }\OtherTok{\textless{}{-}}\NormalTok{ fin\_data }\SpecialCharTok{|\textgreater{}}
  \FunctionTok{filter}\NormalTok{(}\SpecialCharTok{!}\FunctionTok{is.na}\NormalTok{(VIX)) }\SpecialCharTok{|\textgreater{}}
  \FunctionTok{select}\NormalTok{(Asset, Date, Daily\_Return, VIX) }\SpecialCharTok{|\textgreater{}}
  \FunctionTok{collect}\NormalTok{() }\SpecialCharTok{|\textgreater{}}
  \FunctionTok{mutate}\NormalTok{(}\AttributeTok{Date =} \FunctionTok{as.Date}\NormalTok{(Date))}

\NormalTok{q8 }\OtherTok{\textless{}{-}}\NormalTok{ q8\_raw }\SpecialCharTok{|\textgreater{}}
  \FunctionTok{mutate}\NormalTok{(}
    \AttributeTok{VIX\_Regime =} \FunctionTok{case\_when}\NormalTok{(}
\NormalTok{      VIX }\SpecialCharTok{\textgreater{}} \DecValTok{30} \SpecialCharTok{\textasciitilde{}} \StringTok{"High Volatility (VIX \textgreater{} 30)"}\NormalTok{,}
\NormalTok{      VIX }\SpecialCharTok{\textless{}} \DecValTok{20} \SpecialCharTok{\textasciitilde{}} \StringTok{"Low Volatility (VIX \textless{} 20)"}\NormalTok{,}
      \ConstantTok{TRUE} \SpecialCharTok{\textasciitilde{}} \StringTok{"Moderate (20 \textless{}= VIX \textless{}= 30)"}
\NormalTok{    )}
\NormalTok{  ) }\SpecialCharTok{|\textgreater{}}
  \FunctionTok{group\_by}\NormalTok{(Asset, VIX\_Regime) }\SpecialCharTok{|\textgreater{}}
  \FunctionTok{summarise}\NormalTok{(}
    \AttributeTok{Mean\_Daily\_Return =} \FunctionTok{mean}\NormalTok{(Daily\_Return, }\AttributeTok{na.rm =} \ConstantTok{TRUE}\NormalTok{),}
    \AttributeTok{SD\_Daily\_Return =} \FunctionTok{sd}\NormalTok{(Daily\_Return, }\AttributeTok{na.rm =} \ConstantTok{TRUE}\NormalTok{),}
    \AttributeTok{Observations =} \FunctionTok{n}\NormalTok{(),}
    \AttributeTok{.groups =} \StringTok{"drop"}
\NormalTok{  ) }\SpecialCharTok{|\textgreater{}}
  \FunctionTok{arrange}\NormalTok{(Asset, VIX\_Regime)}

\NormalTok{q8\_plot }\OtherTok{\textless{}{-}}\NormalTok{ q8 }\SpecialCharTok{|\textgreater{}}
  \FunctionTok{filter}\NormalTok{(VIX\_Regime }\SpecialCharTok{\%in\%} \FunctionTok{c}\NormalTok{(}\StringTok{"High Volatility (VIX \textgreater{} 30)"}\NormalTok{, }\StringTok{"Low Volatility (VIX \textless{} 20)"}\NormalTok{)) }\SpecialCharTok{|\textgreater{}}
  \FunctionTok{mutate}\NormalTok{(}\AttributeTok{VIX\_Regime =} \FunctionTok{recode}\NormalTok{(VIX\_Regime,}
    \StringTok{"High Volatility (VIX \textgreater{} 30)"} \OtherTok{=} \StringTok{"VIX \textgreater{} 30 (High Fear)"}\NormalTok{,}
    \StringTok{"Low Volatility (VIX \textless{} 20)"} \OtherTok{=} \StringTok{"VIX \textless{} 20 (Low Fear)"}\NormalTok{))}

\FunctionTok{ggplot}\NormalTok{(q8\_plot, }\FunctionTok{aes}\NormalTok{(}\AttributeTok{x =}\NormalTok{ Asset, }\AttributeTok{y =}\NormalTok{ Mean\_Daily\_Return, }\AttributeTok{fill =}\NormalTok{ VIX\_Regime)) }\SpecialCharTok{+}
  \FunctionTok{geom\_col}\NormalTok{(}\AttributeTok{position =} \FunctionTok{position\_dodge}\NormalTok{(}\AttributeTok{width =} \FloatTok{0.7}\NormalTok{), }\AttributeTok{width =} \FloatTok{0.6}\NormalTok{) }\SpecialCharTok{+}
  \FunctionTok{geom\_hline}\NormalTok{(}\AttributeTok{yintercept =} \DecValTok{0}\NormalTok{, }\AttributeTok{linewidth =} \FloatTok{0.3}\NormalTok{, }\AttributeTok{color =} \StringTok{"\#595959"}\NormalTok{) }\SpecialCharTok{+}
  \FunctionTok{scale\_fill\_manual}\NormalTok{(}\AttributeTok{values =} \FunctionTok{c}\NormalTok{(}\StringTok{"VIX \textgreater{} 30 (High Fear)"} \OtherTok{=} \StringTok{"\#C91D42"}\NormalTok{,}
                                \StringTok{"VIX \textless{} 20 (Low Fear)"} \OtherTok{=} \StringTok{"\#2E45B8"}\NormalTok{)) }\SpecialCharTok{+}
  \FunctionTok{scale\_y\_continuous}\NormalTok{(}\AttributeTok{labels =}\NormalTok{ scales}\SpecialCharTok{::}\FunctionTok{percent\_format}\NormalTok{()) }\SpecialCharTok{+}
  \FunctionTok{labs}\NormalTok{(}\AttributeTok{title =} \StringTok{"Mean Daily Return by VIX Regime"}\NormalTok{,}
       \AttributeTok{subtitle =} \StringTok{"Comparing asset performance under high versus low market fear, 2020 to 2025"}\NormalTok{,}
       \AttributeTok{x =} \ConstantTok{NULL}\NormalTok{, }\AttributeTok{y =} \StringTok{"Mean Daily Return"}\NormalTok{, }\AttributeTok{fill =} \ConstantTok{NULL}\NormalTok{) }\SpecialCharTok{+}
  \FunctionTok{theme\_minimal}\NormalTok{(}\AttributeTok{base\_size =} \DecValTok{11}\NormalTok{) }\SpecialCharTok{+}
  \FunctionTok{theme}\NormalTok{(}\AttributeTok{panel.grid.major.x =} \FunctionTok{element\_blank}\NormalTok{(),}
        \AttributeTok{panel.grid.minor =} \FunctionTok{element\_blank}\NormalTok{(),}
        \AttributeTok{axis.ticks =} \FunctionTok{element\_blank}\NormalTok{(),}
        \AttributeTok{plot.title =} \FunctionTok{element\_text}\NormalTok{(}\AttributeTok{face =} \StringTok{"bold"}\NormalTok{, }\AttributeTok{size =} \DecValTok{13}\NormalTok{),}
        \AttributeTok{plot.subtitle =} \FunctionTok{element\_text}\NormalTok{(}\AttributeTok{size =} \DecValTok{9}\NormalTok{, }\AttributeTok{color =} \StringTok{"\#595959"}\NormalTok{),}
        \AttributeTok{legend.position =} \StringTok{"bottom"}\NormalTok{)}
\end{Highlighting}
\end{Shaded}

\begin{table}[H]
\centering
\caption{Q8 Data Output}
\label{tab:q8}
\begin{tabular}{llrrr}
\toprule
Asset & VIX\_Regime & Mean\_Daily\_Return & SD\_Daily\_Return & Obs \\
\midrule
AAPL & High Volatility (VIX > 30) & -0.0034 & 0.0283 & 206 \\
AAPL & Low Volatility (VIX < 20) & 0.0013 & 0.0096 & 1221 \\
AAPL & Moderate (20 <= VIX <= 30) & 0.0004 & 0.0151 & 765 \\
BDO & High Volatility (VIX > 30) & 0.0014 & 0.0263 & 206 \\
BDO & Low Volatility (VIX < 20) & 0.0005 & 0.0131 & 1221 \\
BDO & Moderate (20 <= VIX <= 30) & 0.0011 & 0.0149 & 765 \\
BTC & High Volatility (VIX > 30) & -0.0016 & 0.0495 & 206 \\
BTC & Low Volatility (VIX < 20) & 0.0022 & 0.0261 & 1221 \\
BTC & Moderate (20 <= VIX <= 30) & 0.0016 & 0.0342 & 765 \\
SPY & High Volatility (VIX > 30) & -0.0024 & 0.0220 & 206 \\
SPY & Low Volatility (VIX < 20) & 0.0009 & 0.0052 & 1221 \\
SPY & Moderate (20 <= VIX <= 30) & 0.0002 & 0.0092 & 765 \\
\bottomrule
\end{tabular}
\end{table}


\begin{Shaded}
\begin{Highlighting}[]
\FunctionTok{ggsave}\NormalTok{(}\StringTok{"../reports/figures/fig\_q8.pdf"}\NormalTok{, }\AttributeTok{width =} \DecValTok{7}\NormalTok{, }\AttributeTok{height =} \DecValTok{4}\NormalTok{, }\AttributeTok{device =} \StringTok{"pdf"}\NormalTok{)}
\end{Highlighting}
\end{Shaded}

Table~\ref{tab:q8} categorizes returns conditionally based on the VIX. During high-fear periods (VIX > 30), BTC was the only asset with a positive mean daily return (0.16\%). Every other asset posted negative mean returns. During low-fear periods (VIX < 20), all four assets delivered positive returns, with AAPL leading at 0.25\%.

\subsection{Q9. Interest Rate Sensitivity}

This section addresses the question:

\noindent\textit{How sensitive is each asset's daily return to changes in the 10-year Treasury yield?}

\label{appendix-a9-sql-question-9}




\begin{Shaded}
\begin{Highlighting}[]
\CommentTok{\# Q9: Interest Rate Sensitivity {-}{-} SISON, Aaron Joshua Estacio}

\NormalTok{q9\_raw }\OtherTok{\textless{}{-}}\NormalTok{ fin\_data }\SpecialCharTok{|\textgreater{}}
  \FunctionTok{filter}\NormalTok{(}\SpecialCharTok{!}\FunctionTok{is.na}\NormalTok{(DGS10)) }\SpecialCharTok{|\textgreater{}}
  \FunctionTok{select}\NormalTok{(Asset, Date, Daily\_Return, DGS10) }\SpecialCharTok{|\textgreater{}}
  \FunctionTok{collect}\NormalTok{() }\SpecialCharTok{|\textgreater{}}
  \FunctionTok{mutate}\NormalTok{(}\AttributeTok{Date =} \FunctionTok{as.Date}\NormalTok{(Date))}

\NormalTok{q9 }\OtherTok{\textless{}{-}}\NormalTok{ q9\_raw }\SpecialCharTok{|\textgreater{}}
  \FunctionTok{group\_by}\NormalTok{(Asset) }\SpecialCharTok{|\textgreater{}}
  \FunctionTok{arrange}\NormalTok{(Date) }\SpecialCharTok{|\textgreater{}}
  \FunctionTok{mutate}\NormalTok{(}\AttributeTok{DGS10\_Change =}\NormalTok{ DGS10 }\SpecialCharTok{{-}} \FunctionTok{lag}\NormalTok{(DGS10)) }\SpecialCharTok{|\textgreater{}}
  \FunctionTok{ungroup}\NormalTok{() }\SpecialCharTok{|\textgreater{}}
  \FunctionTok{filter}\NormalTok{(}\SpecialCharTok{!}\FunctionTok{is.na}\NormalTok{(DGS10\_Change)) }\SpecialCharTok{|\textgreater{}}
  \FunctionTok{mutate}\NormalTok{(}
    \AttributeTok{Yield\_Direction =} \FunctionTok{case\_when}\NormalTok{(}
\NormalTok{      DGS10\_Change }\SpecialCharTok{\textgreater{}} \DecValTok{0} \SpecialCharTok{\textasciitilde{}} \StringTok{"Yields Rising"}\NormalTok{,}
\NormalTok{      DGS10\_Change }\SpecialCharTok{\textless{}} \DecValTok{0} \SpecialCharTok{\textasciitilde{}} \StringTok{"Yields Falling"}\NormalTok{,}
      \ConstantTok{TRUE} \SpecialCharTok{\textasciitilde{}} \StringTok{"No Change"}
\NormalTok{    )}
\NormalTok{  ) }\SpecialCharTok{|\textgreater{}}
  \FunctionTok{group\_by}\NormalTok{(Asset, Yield\_Direction) }\SpecialCharTok{|\textgreater{}}
  \FunctionTok{summarise}\NormalTok{(}
    \AttributeTok{Mean\_Daily\_Return =} \FunctionTok{mean}\NormalTok{(Daily\_Return, }\AttributeTok{na.rm =} \ConstantTok{TRUE}\NormalTok{),}
    \AttributeTok{SD\_Daily\_Return =} \FunctionTok{sd}\NormalTok{(Daily\_Return, }\AttributeTok{na.rm =} \ConstantTok{TRUE}\NormalTok{),}
    \AttributeTok{Observations =} \FunctionTok{n}\NormalTok{(),}
    \AttributeTok{.groups =} \StringTok{"drop"}
\NormalTok{  ) }\SpecialCharTok{|\textgreater{}}
  \FunctionTok{arrange}\NormalTok{(Asset, Yield\_Direction)}

\NormalTok{q9\_plot }\OtherTok{\textless{}{-}}\NormalTok{ q9 }\SpecialCharTok{|\textgreater{}}
  \FunctionTok{filter}\NormalTok{(Yield\_Direction }\SpecialCharTok{\%in\%} \FunctionTok{c}\NormalTok{(}\StringTok{"Yields Rising"}\NormalTok{, }\StringTok{"Yields Falling"}\NormalTok{))}

\FunctionTok{ggplot}\NormalTok{(q9\_plot, }\FunctionTok{aes}\NormalTok{(}\AttributeTok{x =}\NormalTok{ Asset, }\AttributeTok{y =}\NormalTok{ Mean\_Daily\_Return, }\AttributeTok{fill =}\NormalTok{ Yield\_Direction)) }\SpecialCharTok{+}
  \FunctionTok{geom\_col}\NormalTok{(}\AttributeTok{position =} \FunctionTok{position\_dodge}\NormalTok{(}\AttributeTok{width =} \FloatTok{0.7}\NormalTok{), }\AttributeTok{width =} \FloatTok{0.6}\NormalTok{) }\SpecialCharTok{+}
  \FunctionTok{geom\_hline}\NormalTok{(}\AttributeTok{yintercept =} \DecValTok{0}\NormalTok{, }\AttributeTok{linewidth =} \FloatTok{0.3}\NormalTok{, }\AttributeTok{color =} \StringTok{"\#595959"}\NormalTok{) }\SpecialCharTok{+}
  \FunctionTok{scale\_fill\_manual}\NormalTok{(}\AttributeTok{values =} \FunctionTok{c}\NormalTok{(}\StringTok{"Yields Rising"} \OtherTok{=} \StringTok{"\#F97A1F"}\NormalTok{,}
                                \StringTok{"Yields Falling"} \OtherTok{=} \StringTok{"\#2E45B8"}\NormalTok{)) }\SpecialCharTok{+}
  \FunctionTok{scale\_y\_continuous}\NormalTok{(}\AttributeTok{labels =}\NormalTok{ scales}\SpecialCharTok{::}\FunctionTok{percent\_format}\NormalTok{()) }\SpecialCharTok{+}
  \FunctionTok{labs}\NormalTok{(}\AttributeTok{title =} \StringTok{"Mean Daily Return by 10{-}Year Treasury Yield Direction"}\NormalTok{,}
       \AttributeTok{subtitle =} \StringTok{"Asset returns conditional on daily yield changes, 2020 to 2025"}\NormalTok{,}
       \AttributeTok{x =} \ConstantTok{NULL}\NormalTok{, }\AttributeTok{y =} \StringTok{"Mean Daily Return"}\NormalTok{, }\AttributeTok{fill =} \ConstantTok{NULL}\NormalTok{) }\SpecialCharTok{+}
  \FunctionTok{theme\_minimal}\NormalTok{(}\AttributeTok{base\_size =} \DecValTok{11}\NormalTok{) }\SpecialCharTok{+}
  \FunctionTok{theme}\NormalTok{(}\AttributeTok{panel.grid.major.x =} \FunctionTok{element\_blank}\NormalTok{(),}
        \AttributeTok{panel.grid.minor =} \FunctionTok{element\_blank}\NormalTok{(),}
        \AttributeTok{axis.ticks =} \FunctionTok{element\_blank}\NormalTok{(),}
        \AttributeTok{plot.title =} \FunctionTok{element\_text}\NormalTok{(}\AttributeTok{face =} \StringTok{"bold"}\NormalTok{, }\AttributeTok{size =} \DecValTok{13}\NormalTok{),}
        \AttributeTok{plot.subtitle =} \FunctionTok{element\_text}\NormalTok{(}\AttributeTok{size =} \DecValTok{9}\NormalTok{, }\AttributeTok{color =} \StringTok{"\#595959"}\NormalTok{),}
        \AttributeTok{legend.position =} \StringTok{"bottom"}\NormalTok{)}
\end{Highlighting}
\end{Shaded}

\begin{table}[H]
\centering
\caption{Q9 Data Output}
\label{tab:q9}
\begin{tabular}{llrrr}
\toprule
Asset & Yield\_Direction & Mean\_Daily\_Return & SD\_Daily\_Return & Obs \\
\midrule
AAPL & No Change      & 0.0019 & 0.0195 & 103 \\
AAPL & Yields Falling & 0.0005 & 0.0200 & 684 \\
AAPL & Yields Rising  & 0.0013 & 0.0201 & 712 \\
BDO  & No Change      & -0.0020 & 0.0208 & 103 \\
BDO  & Yields Falling & 0.0000 & 0.0253 & 684 \\
BDO  & Yields Rising  & 0.0007 & 0.0201 & 712 \\
BTC  & No Change      & 0.0028 & 0.0372 & 103 \\
BTC  & Yields Falling & 0.0017 & 0.0343 & 684 \\
BTC  & Yields Rising  & 0.0020 & 0.0363 & 712 \\
SPY  & No Change      & -0.0001 & 0.0101 & 103 \\
SPY  & Yields Falling & -0.0001 & 0.0130 & 684 \\
SPY  & Yields Rising  & 0.0014 & 0.0135 & 712 \\
\bottomrule
\end{tabular}
\end{table}

\begin{Shaded}
\begin{Highlighting}[]
\FunctionTok{ggsave}\NormalTok{(}\StringTok{"../reports/figures/fig\_q9.pdf"}\NormalTok{, }\AttributeTok{width =} \DecValTok{7}\NormalTok{, }\AttributeTok{height =} \DecValTok{4}\NormalTok{, }\AttributeTok{device =} \StringTok{"pdf"}\NormalTok{)}
\end{Highlighting}
\end{Shaded}

Table~\ref{tab:q9} provides the conditional means measured against the directional movement of the 10-year Treasury yield. During periods when the yield rose, all four assets delivered positive mean daily returns. During falling-yield regimes, SPY recorded a mean return of -0.01\%, while BDO remained largely insensitive at 0.00\%.

\subsection{Q10. Inflation Hedge Performance}

This section addresses the question:

\noindent\textit{Which assets served as the most effective inflation hedges during the high-inflation period of 2021 to 2023?}

\label{appendix-a10-sql-question-10}




\begin{Shaded}
\begin{Highlighting}[]
\CommentTok{\# Q10: Inflation Hedge Performance {-}{-} SISON, Aaron Joshua Estacio}

\NormalTok{q10\_raw }\OtherTok{\textless{}{-}}\NormalTok{ fin\_data }\SpecialCharTok{|\textgreater{}}
  \FunctionTok{select}\NormalTok{(Date, Asset, Daily\_Return, CPI\_YoY) }\SpecialCharTok{|\textgreater{}}
  \FunctionTok{collect}\NormalTok{() }\SpecialCharTok{|\textgreater{}}
  \FunctionTok{mutate}\NormalTok{(}\AttributeTok{Date =} \FunctionTok{as.Date}\NormalTok{(Date))}

\NormalTok{q10 }\OtherTok{\textless{}{-}}\NormalTok{ q10\_raw }\SpecialCharTok{|\textgreater{}}
  \FunctionTok{filter}\NormalTok{(}\SpecialCharTok{!}\FunctionTok{is.na}\NormalTok{(CPI\_YoY) }\SpecialCharTok{\&}\NormalTok{ CPI\_YoY }\SpecialCharTok{\textgreater{}} \FloatTok{0.05}\NormalTok{) }\SpecialCharTok{|\textgreater{}}
  \FunctionTok{filter}\NormalTok{(}\SpecialCharTok{!}\FunctionTok{is.na}\NormalTok{(Daily\_Return)) }\SpecialCharTok{|\textgreater{}}
  \FunctionTok{group\_by}\NormalTok{(Asset) }\SpecialCharTok{|\textgreater{}}
  \FunctionTok{arrange}\NormalTok{(Date) }\SpecialCharTok{|\textgreater{}}
  \FunctionTok{mutate}\NormalTok{(}\AttributeTok{Cumulative\_Return =} \FunctionTok{cumprod}\NormalTok{(}\DecValTok{1} \SpecialCharTok{+}\NormalTok{ Daily\_Return) }\SpecialCharTok{{-}} \DecValTok{1}\NormalTok{) }\SpecialCharTok{|\textgreater{}}
  \FunctionTok{summarise}\NormalTok{(}
    \AttributeTok{Start\_Date =} \FunctionTok{min}\NormalTok{(Date),}
    \AttributeTok{End\_Date =} \FunctionTok{max}\NormalTok{(Date),}
    \AttributeTok{Start\_CPI\_YoY =} \FunctionTok{first}\NormalTok{(CPI\_YoY),}
    \AttributeTok{End\_CPI\_YoY =} \FunctionTok{last}\NormalTok{(CPI\_YoY),}
    \AttributeTok{Total\_Cumulative\_Return =} \FunctionTok{last}\NormalTok{(Cumulative\_Return),}
    \AttributeTok{Mean\_Daily\_Return =} \FunctionTok{mean}\NormalTok{(Daily\_Return, }\AttributeTok{na.rm =} \ConstantTok{TRUE}\NormalTok{),}
    \AttributeTok{Observations =} \FunctionTok{n}\NormalTok{()}
\NormalTok{  ) }\SpecialCharTok{|\textgreater{}}
  \FunctionTok{arrange}\NormalTok{(}\FunctionTok{desc}\NormalTok{(Total\_Cumulative\_Return))}

\NormalTok{q10\_plot }\OtherTok{\textless{}{-}}\NormalTok{ q10 }\SpecialCharTok{|\textgreater{}}
  \FunctionTok{mutate}\NormalTok{(}\AttributeTok{Asset =} \FunctionTok{factor}\NormalTok{(Asset, }\AttributeTok{levels =}\NormalTok{ q10}\SpecialCharTok{$}\NormalTok{Asset[}\FunctionTok{order}\NormalTok{(q10}\SpecialCharTok{$}\NormalTok{Total\_Cumulative\_Return)]))}

\FunctionTok{ggplot}\NormalTok{(q10\_plot, }\FunctionTok{aes}\NormalTok{(}\AttributeTok{x =}\NormalTok{ Total\_Cumulative\_Return, }\AttributeTok{y =}\NormalTok{ Asset, }\AttributeTok{fill =}\NormalTok{ Asset)) }\SpecialCharTok{+}
  \FunctionTok{geom\_col}\NormalTok{(}\AttributeTok{width =} \FloatTok{0.6}\NormalTok{, }\AttributeTok{show.legend =} \ConstantTok{FALSE}\NormalTok{) }\SpecialCharTok{+}
  \FunctionTok{geom\_text}\NormalTok{(}\FunctionTok{aes}\NormalTok{(}\AttributeTok{x =}\NormalTok{ Total\_Cumulative\_Return }\SpecialCharTok{/} \DecValTok{2}\NormalTok{,}
                \AttributeTok{label =} \FunctionTok{sprintf}\NormalTok{(}\StringTok{"\%+.1f\%\%"}\NormalTok{, Total\_Cumulative\_Return }\SpecialCharTok{*} \DecValTok{100}\NormalTok{)),}
            \AttributeTok{hjust =} \FloatTok{0.5}\NormalTok{, }\AttributeTok{size =} \FloatTok{3.2}\NormalTok{, }\AttributeTok{color =} \StringTok{"white"}\NormalTok{, }\AttributeTok{fontface =} \StringTok{"bold"}\NormalTok{) }\SpecialCharTok{+}
  \FunctionTok{scale\_fill\_manual}\NormalTok{(}\AttributeTok{values =} \FunctionTok{c}\NormalTok{(}\AttributeTok{AAPL =} \StringTok{"\#555555"}\NormalTok{, }\AttributeTok{BDO =} \StringTok{"\#003D6B"}\NormalTok{,}
                                \AttributeTok{BTC =} \StringTok{"\#F7931A"}\NormalTok{, }\AttributeTok{SPY =} \StringTok{"\#00733E"}\NormalTok{)) }\SpecialCharTok{+}
  \FunctionTok{scale\_x\_continuous}\NormalTok{(}\AttributeTok{labels =}\NormalTok{ scales}\SpecialCharTok{::}\FunctionTok{percent\_format}\NormalTok{(),}
                     \AttributeTok{expand =} \FunctionTok{expansion}\NormalTok{(}\AttributeTok{mult =} \FunctionTok{c}\NormalTok{(}\FloatTok{0.15}\NormalTok{, }\FloatTok{0.15}\NormalTok{))) }\SpecialCharTok{+}
  \FunctionTok{labs}\NormalTok{(}\AttributeTok{title =} \StringTok{"Cumulative Return During High Inflation Period"}\NormalTok{,}
       \AttributeTok{subtitle =} \StringTok{"June 2021 to February 2023, CPI year{-}on{-}year above 5\%"}\NormalTok{,}
       \AttributeTok{x =} \StringTok{"Cumulative Return"}\NormalTok{, }\AttributeTok{y =} \ConstantTok{NULL}\NormalTok{) }\SpecialCharTok{+}
  \FunctionTok{theme\_minimal}\NormalTok{(}\AttributeTok{base\_size =} \DecValTok{11}\NormalTok{) }\SpecialCharTok{+}
  \FunctionTok{theme}\NormalTok{(}\AttributeTok{panel.grid.major.y =} \FunctionTok{element\_blank}\NormalTok{(),}
        \AttributeTok{panel.grid.minor =} \FunctionTok{element\_blank}\NormalTok{(),}
        \AttributeTok{axis.ticks =} \FunctionTok{element\_blank}\NormalTok{(),}
        \AttributeTok{plot.title =} \FunctionTok{element\_text}\NormalTok{(}\AttributeTok{face =} \StringTok{"bold"}\NormalTok{, }\AttributeTok{size =} \DecValTok{13}\NormalTok{),}
        \AttributeTok{plot.subtitle =} \FunctionTok{element\_text}\NormalTok{(}\AttributeTok{size =} \DecValTok{9}\NormalTok{, }\AttributeTok{color =} \StringTok{"\#595959"}\NormalTok{))}
\end{Highlighting}
\end{Shaded}

\begin{table}[H]
\centering
\caption{Q10 Data Output}
\label{tab:q10}
\begin{tabular}{lllrrrrr}
\toprule
Asset & Start\_Date & End\_Date & Start\_CPI\_YoY & End\_CPI\_YoY & Total\_Cumulative\_Return & Mean\_Daily\_Return & Obs \\
\midrule
BDO & 2021-06-01 & 2023-02-28 & 0.0530 & 0.0596 & 0.4531 & 0.0010 & 456 \\
AAPL & 2021-06-01 & 2023-02-28 & 0.0530 & 0.0596 & 0.2865 & 0.0007 & 456 \\
SPY & 2021-06-01 & 2023-02-28 & 0.0530 & 0.0596 & 0.0243 & 0.0001 & 456 \\
BTC & 2021-06-01 & 2023-02-28 & 0.0530 & 0.0596 & -0.3800 & -0.0002 & 456 \\
\bottomrule
\end{tabular}
\end{table}


\begin{Shaded}
\begin{Highlighting}[]
\FunctionTok{ggsave}\NormalTok{(}\StringTok{"../reports/figures/fig\_q10.pdf"}\NormalTok{, }\AttributeTok{width =} \DecValTok{7}\NormalTok{, }\AttributeTok{height =} \FloatTok{3.5}\NormalTok{, }\AttributeTok{device =} \StringTok{"pdf"}\NormalTok{)}
\end{Highlighting}
\end{Shaded}

Table~\ref{tab:q10} records the cumulative returns during the high-inflation window of June 2021 to February 2023. BDO generated a cumulative return of 43.9\%, AAPL gained 19.5\%, SPY lost 3.2\%, and BTC collapsed by 38.0\%. The data query directly challenges prevailing market narratives regarding digital inflation hedges.
